% Sixfold contribution, GPT-6 Astra Pro /root/sixfolds, 2026-11-21.
% The source proofs below were read in full; exact finite checks are in
% sixfold_certificates.py and sixfold_division_certificates.py.
\section{Sixfolds: the exact cyclic source and explicit bounds}
\label{sec:sixfolds}

The natural genus-four cyclic construction reaches precisely the two
quadratic cyclotomic fields and the split rational Hermitian class.
The cubic and quartic period maps both have differential rank nine.
The published dominance results are cited at their exact locators below;
the character, determinant, degree, lattice and division consequences are
proved separately. Throughout, determinant means its class in
$\mathbf Q^\times/\operatorname{Nm}_{K/\mathbf Q}(K^\times)$,
without absorbing a dimension-dependent sign.

\begin{proposition}[SF-01: the cyclic character-field restriction]
\label{SF-01}
Let $C\to X$ be a connected etale cyclic cover with group $C_d$ and
$g(X)=4$. For $e>1$ dividing $d$, let $B_e$ be its rational cyclotomic
isotypic abelian factor of order $e$. Then
\[
 \dim B_e=3\varphi(e).
\]
If $B_e$ is a sixfold, then $e=3,4,$ or $6$, and its coefficient
field is respectively $\mathbf Q(\sqrt{-3}),\mathbf Q(i),$ or
$\mathbf Q(\sqrt{-3})$.
\end{proposition}
\begin{proof}
Write $L$ for the order-$d$ line bundle defining the cover. The nontrivial
holomorphic eigenspaces are $H^0(X,\omega_X\otimes L^j)$. Riemann--Roch
and $H^0(X,L^{-j})=0$ give dimension $g(X)-1=3$. There are
$\varphi(e)$ characters of exact order $e$, giving the displayed
formula. The equation $\varphi(e)=2$ has exactly the three stated
solutions. Passing from the original cover to its quotient of degree
$e$ does not change this cyclotomic constituent.
\end{proof}
This is a restriction on the prescribed rational isotypic construction.
It does not deny that an individual special factor could have accidental
additional endomorphisms. Such special points do not establish dominance
in a general Weil family with generic endomorphism algebra $K$.

\begin{proposition}[SF-02: the cyclic Hermitian space is split]
\label{SF-02}
For a connected etale cyclic cover of a genus-$(n+1)$ curve, the
order-$e$ primitive homology in the quadratic cases $e=3,4,6$ is a
split Hermitian space of $K$-rank $2n$. Its determinant is
$[(-1)^n]$. In particular every sixfold in SF-01 belongs to determinant
class $[-1]$.
\end{proposition}
\begin{proof}
A connected cover is specified by a surjection
$\chi:H_1(X,\mathbf Z)\to\mathbf Z/d$. It has a primitive integral
lift. Here is a direct justification. Choose integer representatives
$v_1,\ldots,v_{2g}$ and put $b=\gcd(v_2,\ldots,v_{2g})$.
If $b=0$, replace one zero coordinate by $d$. For primes dividing both
$b$ and $d$, surjectivity implies that $v_1$ is not divisible by that
prime. For each prime dividing $b$ but not $d$, choose $k$ so that
$v_1+kd$ is nonzero modulo that prime. The Chinese remainder theorem
makes these finitely many choices simultaneously. The resulting lift
has relatively prime coordinates. A primitive vector in a unimodular
symplectic lattice extends to a symplectic basis. Surjectivity of the
surface mapping-class representation onto $\operatorname{Sp}_{2g}(\mathbf Z)$
\cite[Theorem~6.4]{SF-FarbMargalit} therefore puts all monodromy on the
first handle.

Decompose $X$ into a one-holed torus and a genus-$n$ surface with one
boundary circle. The latter has $d$ disjoint lifts; the connected cover
of the former has genus one and $d$ boundary circles. On each unaffected
handle choose cycles $a_{j,k},b_{j,k}$, $1\le j\le n$,
$k\in\mathbf Z/d$, with intersection form
\[
 (a_{j,k},a_{l,r})=(b_{j,k},b_{l,r})=0,
 \qquad (a_{j,k},b_{l,r})=\delta_{jl}\delta_{kr}.
\]
The deck generator shifts the sheet index. These cycles give a
nondegenerate rank-$2nd$ summand. Mayer--Vietoris shows that the
complement is the rank-two deck-invariant homology of the covered
torus after its boundary circles are filled by the other pieces.
Consequently all nontrivial isotypic homology comes from
\[
 \mathbf Q[C_d]^n\oplus\mathbf Q[C_d]^n,
\]
whose two displayed summands are invariant, totally isotropic and
paired perfectly. The order-$e$ idempotent yields two $K$-subspaces of
dimension $n$. They are Hermitian-isotropic: all field multiples of a
vector stay in the corresponding symplectically isotropic subspace.
Nondegeneracy now gives a Hermitian matrix
\[
 \begin{pmatrix}0&A\\ A^*&0\end{pmatrix},
 \qquad \det=(-1)^n\operatorname{Nm}(\det A).
\]
This proves the assertion.
\end{proof}

\paragraph{The missing components.}
The union of the cubic and quartic sources has image dimension exactly
nine in the split Eisenstein and Gaussian families, by the rank checks
below. It has no image in the other determinant components as a map
preserving the prescribed $K$-action and polarization. An empty image
is not a period map of rank zero. Among the class-number-one fields,
all components for
\[
 \mathbf Q(\sqrt{-2}),\ \mathbf Q(\sqrt{-7}),\
 \mathbf Q(\sqrt{-11}),\ \mathbf Q(\sqrt{-19}),\
 \mathbf Q(\sqrt{-43}),\ \mathbf Q(\sqrt{-67}),\
 \mathbf Q(\sqrt{-163})
\]
are absent, as are all nonsplit Gaussian and Eisenstein components.
The standard class-number-one list is used only to spell out these
names \cite[``Class Number 1'']{SF-Stein}; the propositions themselves
have no class-number hypothesis. Explicit missing Hermitian forms are
\[
 \operatorname{diag}(1,1,1,-1,-1,-3)\quad(K=\mathbf Q(i)),
 \qquad
 \operatorname{diag}(1,1,1,-1,-1,-2)\quad(K=\mathbf Q(\sqrt{-3})).
\]
Their determinants are $-3$ and $-2$. The positive factors $3$ and $2$
are not norms in the respective fields: the primes are inert, so norm
valuations are even. Both signatures are $(3,3)$. These are
counterexamples to universal cyclic-source dominance, not to the Hodge
conjecture.

\subsection{Published differential ranks with exact finite checks}
\label{sec:sixfold-ranks}

\paragraph{SF-CW: the cubic rank is nine.}
Schoen \cite[Lemma~13, equations~(13.1)--(13.2), pp.~334--335]{SF-Schoen1998}
proves the dominance of the cubic sixfold source. His cover of
$\mathbf P^1$ has group $(\mathbf Z/3)^3$, inertia
$e_1,e_2,e_3,-e_1-e_2-e_3$, and etale subgroup generated by $e_1+e_2$.
The codifferential is multiplication of the two three-dimensional
character spaces. Their one-dimensional constituent characters are
\[
 \begin{split}
 U_+ &: (0,1,1),\ (1,0,1),\ (2,2,1),\\
 U_- &: (1,1,0),\ (1,1,1),\ (1,1,2).
 \end{split}
\]
The nine sums modulo three are
\[
 \begin{array}{ccc}
 (1,2,1)&(1,2,2)&(1,2,0)\\
 (2,1,1)&(2,1,2)&(2,1,0)\\
 (0,0,1)&(0,0,2)&(0,0,0).
 \end{array}
\]
They are distinct. Each product is nonzero because the curve is smooth
and integral. Thus the nine products are independent and the
codifferential has exact rank nine. Source and target both have
dimension nine. The Chevalley--Weil multiplicities and this finite
character computation were independently recomputed in
\texttt{sixfold\_certificates.py}; the dominance theorem remains a
published cited input.

\paragraph{SF-KO: the quartic rank is nine.}
Koike \cite[Theorem~2.1 and Proposition~3.1]{SF-Koike2004} resolves the
sixfold question recorded in van Geemen
\cite[Remark~5.4, p.~630]{SF-vanGeemen1996}. Koike's codifferential is
\[
 H^0(X,\omega_XL)\otimes H^0(X,\omega_XL^3)
 \longrightarrow H^0(X,\omega_X^2),\qquad 3\cdot3\longrightarrow9.
\]
Its exact locator is equation~(1), p.~568; the eigenvectors are (2),
p.~568; the relations are (6)--(7), p.~571; coefficient substitution
(8) and monomial basis (9) occur on pp.~571--572. The source consists
of invariant complete intersections in $(\mathbf P^1)^4$ with free
cyclic-four action. Three normalized product coordinates give an
identity block. In the remaining six coordinates, the polynomial
matrix $M$ specializes formally at $A=B=C=0$ to
\[
 M_0=\begin{pmatrix}
 1&0&0&0&0&0\\
 0&1&0&0&-2&0\\
 0&0&1&0&0&0\\
 0&0&0&1&0&0\\
 0&0&0&0&-1&0\\
 0&0&0&0&0&-1
 \end{pmatrix},\qquad \det M_0=1.
\]
This point is a witness that the determinant polynomial is nonzero;
it is not asserted to be a smooth curve. The nonempty rank-open set
therefore intersects the nonempty smooth/free-action open set of
Proposition~3.1 in the irreducible coefficient space. The differential
rank at such smooth points is nine.

As a separate finite check, at $A=(1,2,3)$, $B=(2,3,5)$,
$C=(3,5,7)$ the six-product determinant is $-389639$, and the minor
formed by the first six rows of the seven quadratic $b$-relations is
$227264$. All sixteen polynomial identities in (6)--(7) were expanded
exactly and verified over $\mathbf Q(i)$ in
\texttt{sixfold\_certificates.py}. This second parameter point also
serves only as a rank-polynomial witness. Koike's printed symbol
$\delta=1$ is not silently identified with our unsigned determinant
class; SF-02 establishes our convention $[-1]$.

\subsection{Integral detector degrees and the polarization}

\begin{theorem}[SF-03: the Gaussian principal Prym sixfold bound]
\label{SF-03}
Let $C\to X$ be an etale cyclic-four cover with $g(X)=4$, deck
generator $\sigma$, and
\[
 B=\operatorname{Prym}(C\to C/\langle\sigma^2\rangle),\qquad
 P=\ker(\operatorname{Nm}_{C/X})^0.
\]
Write $\Theta=\Theta_C$ and $\Theta|_B=2\Xi$, where $\Xi$ is the
integral principal Prym polarization. Then
\[
 \dim J(C)=13,\qquad \dim P=9,\qquad\dim B=6,
\]
$\Theta|_P$ has type $(1^6,4^3)$, and an integral effective detecting
threefold in $B$ has $\Xi$-degree at most
\[
 D_{G,6}=121\,993\,819\,950\,101\,299\,200.
\]
The same bound holds on every fibre of the connected principal Prym
Weil component dominated by these covers.
\end{theorem}
\begin{proof}
Use the unramified, $r=0$, supplier construction in Schoen
\cite[Theorem~2.0, pp.~12--13, Lemma~2.6, and \S3, pp.~24--25]{SF-Schoen1988}.
For $h=6$ its supplier cycles on $C^6$ are supported over
$|K_X|\subset\operatorname{Sym}^6X$. We spell out the detecting
component argument because its support determines the degree bound.

Let $a:C^6\to J(C)$ be the sum Abel map and $p=1-\sigma^2:J(C)\to B$.
The pullback $(p\circ a)^*$ maps the two Weil lines injectively to the
two alternating top tensors of the primitive $\pm i$ eigenspaces in
$H^1(C)^{\otimes6}$. Products of two forms from the same such
eigenspace on one curve factor vanish: their character is nontrivial
whereas $H^2(C)$ has trivial character. Thus the determinant pullback
has exactly this primitive Kunneth component. Its middle pairing is
nondegenerate, since the pairing of the two eigenspaces is perfect
and the pairing of their top exterior powers is its nonzero
determinant. The relevant primitive Kunneth projections of Schoen's
supplier classes span this summand; other Kunneth components pair
trivially with it. Therefore one supplier has nonzero pairing with a
nonzero pulled-back Weil vector. Upon decomposing it into integral
irreducible components, at least one effective component still has
nonzero pairing.

Its Abel image $Z$ and its image $p(Z)$ both have dimension three;
otherwise the corresponding pushforward would vanish in that
dimension. Each proper pushforward is a positive integral multiple
of the reduced image, so reduced images preserve detection. The norm
of the supplier divisor has class $K_X$. Hence
\[
 Z\subset W_6(C)\cap T_0,
\]
where $T_0$ is a connected component of the appropriate norm fibre,
a translate of the full $P$, not of $B$. Conversely an effective
divisor representing a point of this intersection has norm in
$|K_X|$. The map $\operatorname{Sym}^6C\to\operatorname{Sym}^6X$ is
finite. Its inverse image of the three-dimensional canonical system
has dimension at most three. It follows that
$\dim(W_6(C)\cap T_0)\le3$, and $Z$ is an irreducible component of
that intersection.

The standard etale cyclic Prym type gives $\Theta|_P$ type
$(1^6,4^3)$ and Euler characteristic $64$; see
\cite[\S4, polarization-type paragraph, author version p.~6]{SF-LangeOrtega2011},
which cites Birkenhake--Lange Corollary~12.1.4 and Lemma~12.3.1.
Equivalently, the complementary factor $\pi^*J(X)$ has induced type
$(1,4,4,4)$: the degree-four isogeny from $J(X)$ has cyclic kernel
and pulls the restricted form back to $4\Theta_X$. Complementarity
in the principal $J(C)$ gives the asserted type of $P$.
The intermediate double cover is etale, so the double-Prym theorem
gives $\Theta|_B=2\Xi$ with $\Xi$ principal
\cite[\S3, Corollary~2 after Corollary~1, p.~332]{SF-Mumford1974}.

In the embedding by $3\Theta$, Poincare's formula gives
\[
 \deg_{3\Theta} W_6(C)=3^6\frac{13!}{7!}=900\,694\,080,
 \qquad
 \deg_{3\Theta} T_0=3^9\,9!\,64=457\,124\,290\,560.
\]
The component Bezout bound, valid for improper intersections
\cite[Example~8.4.6]{SF-Fulton}, therefore gives
\[
 \deg_\Theta Z\le
 \frac{900\,694\,080\cdot457\,124\,290\,560}{27}.
\]
The deck action preserves $\Theta$ and has Rosati adjoint
$\sigma^{-1}$. Thus
\[
 p^\dagger p=2(1-\sigma^2)\le4\operatorname{id},
 \quad p^*(\Theta|_B)\le4\Theta,
 \quad p^*\Xi\le2\Theta
\]
in the semipositive order on invariant forms. The projection formula
for the effective threefold gives
\[
 \deg_\Xi(p_*Z)\le8\deg_\Theta Z.
\]
Removing its positive generic-degree multiplicity only decreases the
degree and leaves a detector. The resulting integer is
\[
 \frac8{27}\cdot900\,694\,080\cdot457\,124\,290\,560
 =121\,993\,819\,950\,101\,299\,200.
\]
No rational division of a cycle was used. Koike's rank-nine dominance
supplies a dense set of fibres; the authorized detector closedness
input, Theorem~D/Lemma~Dom of the division attack, extends this same
integer bound to the whole component.
\end{proof}

\paragraph{Why the full norm fibre matters.}
For a cyclic-four cover of a genus-$(n+1)$ curve the supplier norm
fibre has dimension $3n$ and polarization Euler characteristic $4^n$.
The primitive double Prym has dimension $2n$ and Euler characteristic
$2^{2n}=4^n$. The equal Euler characteristics do not permit replacing
the norm fibre's factorial $(3n)!$ and power $3^{3n}$ by those of
the smaller primitive Prym.

\paragraph{The Eisenstein source normalization.}
The supplied GHOSTS R20--R24 input uses the cubic cover
$C_{10}\to X_4$ and the full six-dimensional norm kernel $P_6$,
with restricted Jacobian polarization of type $(1^3,3^3)$. Its source
detector bound is
\[
 D_{E,6}=57\,854\,793\,024\,000
 =\frac{110\,224\,800\cdot14\,171\,760}{27},
\]
where
\[
 110\,224\,800=3^6\frac{10!}{4!},\qquad
 14\,171\,760=3^6\,6!\,27.
\]
The supplier already lies in a translate of this full $P_6$, so no
$1-\sigma$ projection cost occurs in that calculation. This
normalization is not called principal. An isogenous principal
quotient or another polarized lattice requires its own charged
lattice map. The source statement is used as authorized background;
its arithmetic was recomputed, but the earlier full proof is not
present in the supplied archive and is not claimed independently
reproved here.

\begin{proposition}[SF-04: a fixed split-polarization transfer]
\label{SF-04}
Fix a polarized Weil sixfold lattice with rational Hermitian space
isometric to the Gaussian source of SF-03 or the cubic source just
described. Choose one rational $K$-linear isometry
$F:V_{\rm source}\to V_{\rm target}$ of their alternating
polarization forms, and let $m\ge1$ clear its integral-lattice
denominators. Then on the target component
\[
 \delta_W\le m^6D_{\rm source},
 \qquad D_{\rm source}=D_{G,6}\ \hbox{or}\ D_{E,6}.
\]
\end{proposition}
\begin{proof}
The isometry identifies the period domains. Intersecting the two
lattice-stabilizer groups gives a finite-index subgroup, so after a
finite level cover $mF$ induces a family of $K$-equivariant
isogenies $f:B\to A$ with
$f^*c_1(L)=m^2c_1(L_{\rm source})$. For an effective source detector
$Y$ of dimension three,
\[
 \deg_L(f_*Y)=\int_Y f^*c_1(L)^3
 =m^6\deg_{L_{\rm source}}Y.
\]
On the rational Weil plane $f_*$ is an isomorphism and respects the
$K$-character decomposition. The reduced image is still detecting
and has no larger degree. Source dominance and detector closedness
give the bound throughout the target component. Existence of $F$
for signature $(3,3)$ and determinant $[-1]$ follows from rational
Hermitian classification, Schoen \cite[\S7, p.~332]{SF-Schoen1998};
the family isogeny is explained in \S6 on that page. Enumerating
rational $K$-matrices and testing $F^*H_{\rm target}F=H_{\rm source}$
terminates by this classification, after which $m$ is an exact
integer.
\end{proof}
The isometry and denominator are fixed for one polarized family.
This does not assert control by a varying isogeny degree or the core
index $J$ on a Hecke orbit. Rescaling a rank-six Hermitian form by a
rational $c$ multiplies its determinant by $c^6=\operatorname{Nm}(c^3)$;
it cannot change its determinant class.

\begin{proposition}[SF-05: both maximal-order Gaussian principal types]
\label{SF-05}
For an integral maximal $\mathbf Z[i]$ action compatible with a
principal polarization, a balanced rank-$2n$ Gaussian Hermitian
lattice has exactly two possible isometry types:
\[
 U^n,\quad U=\begin{pmatrix}0&1\\1&0\end{pmatrix},
 \qquad I_{n,n}=\operatorname{diag}(1^n,-1^n).
\]
The cyclic-four principal Prym has the first type. A fixed isogeny
from the first type to the second pulls back the principal
polarization to twice the source principal polarization. Thus the
sixfold bound valid on either maximal-order principal type is
\[
 D_{G,6}^{\rm principal}=8D_{G,6}
 =975\,950\,559\,600\,810\,393\,600.
\]
\end{proposition}
\begin{proof}
The Gaussian module is free. If $E$ is the principal alternating
form, its associated integral Hermitian matrix satisfies
$\det_{\mathbf Z}E=(\det_K H)^2$. Hence $H$ is unimodular and its
signature gives determinant $(-1)^n$. Rational Hermitian
classification makes its rational space split. Choose a rational
isotropic vector and divide an integral multiple by its Gaussian
content to obtain a primitive isotropic lattice vector $e$.
Unimodularity gives $f$ with $H(e,f)=1$. The span has matrix
\[
 \begin{pmatrix}0&1\\1&k\end{pmatrix},\quad k\in\mathbf Z,
\]
and determinant $-1$, hence splits integrally. Replacing $f$ by
$f-ce$ changes $k$ by $2\operatorname{Re}(c)$. An even $k$ reduces
to $U$; an odd $k$ reduces to the block with $k=1$, whose basis
$f,e-f$ has matrix $I_{1,1}$. The complement is balanced,
unimodular and rationally split, so induction gives a sum of these
two sorts of planes. If an odd block occurs, convert every remaining
$U$ using
\[
 \langle1\rangle\oplus U\simeq\langle1,1,-1\rangle:
 \quad a+e,\ a-f,\ a+e-f
\]
is an orthogonal unimodular basis with norms $1,1,-1$ when
$(a,e,f)$ has matrix $\operatorname{diag}(1,U)$. Retaining the
negative line of the original odd block gives $I_{2,2}$, and the
process continues. Parity of vector norms distinguishes the even
$U^n$ from the odd $I_{n,n}$.

For the source, the topological model of SF-02 gives on every
unaffected handle the saturated primitive integral basis
\[
 a_0-a_2,\quad a_1-a_3,\quad b_0-b_2,\quad b_1-b_3.
\]
Its Jacobian intersection form is twice symplectic and its half is
$\Xi$. The deck action takes the first difference to the second and
the second to minus the first. In the Gaussian basis
$a_0-a_2,b_0-b_2$ the Hermitian form is
$\bigl(\begin{smallmatrix}0&i\\-i&0\end{smallmatrix}\bigr)$,
which becomes $U$ by a unit change. The covered torus contributes
only invariant homology, so these blocks exhaust the primitive
lattice. The published lattice calculation is also in van Geemen
\cite[Theorem~5.3(1)--(2), p.~629]{SF-vanGeemen1996}.

Finally, with
$M=\bigl(\begin{smallmatrix}1&1\\1&-1\end{smallmatrix}\bigr)$,
\[
 M^*I_{1,1}M=2U.
\]
Direct sums give an integral Gaussian map, invertible over $K$,
which transports the complex structure and identifies the two
period domains. Its induced isogeny has degree $2^{2n}$ and
pullback polarization $2\Xi$. The degree of a dimension-$n$
effective cycle grows by $2^n$ by the projection formula. Detection
is preserved by the $K$-isogeny and reduced images lower degree.
Putting $n=3$ gives the stated factor eight.
\end{proof}
For a rational action or a nonmaximal integral order, the two-lattice
conclusion is not asserted; SF-04 retains the fixed denominator cost.

\subsection{Small integral projectors and quadratic division constants}

\begin{theorem}[SF-06: an explicit Gaussian sixfold division constant]
\label{SF-06}
Let a polarized Gaussian Weil sixfold have integral $\mathbf Z[i]$
action and an effective threefold detector of degree at most $D$.
For $0\ne t\in W$, with $\delta t$ integral and
$T=-\int t^2>0$, one has
\[
 c_t\le\delta D^2\max\{31\,719\,424,
              563\,200\,000\lceil\sqrt T\rceil\}.
\]
Here degrees use the stated integral compatible polarization;
principality is not needed for this division step.
\end{theorem}
\begin{proof}
On $H^6$ put $I=[i]^*$, $A=[2+i]^*$ and
$\bar A=[2-i]^*$. In the character decomposition indexed by
$j=0,\ldots,6$, the eigenvalue of $I$ is $(-1)^{j+1}$.
On the even indices, the eigenvalues of $A+\bar A$ are $-234$
for $j=0,6$ and $150$ for $j=2,4$. Therefore the integral
correspondence
\[
 (1-I)(150-A-\bar A)=768e_W.
\]
For a detector $Y$ set $y=768e_W[Y]$. The expansion
\[
 y=(150+IA+I\bar A)[Y]-(150I+A+\bar A)[Y]
\]
has integral effective positive and negative parts, each of degree
at most $(150+125+125)D=400D$, because the pullbacks by $2\pm i$
multiply threefold degree by $5^3=125$. The authorized orthogonal
projector and comass estimates of Theorem~W give
\[
 d=-\int y^2\in\mathbf Z_{>0},\qquad
 \sqrt d\le768D/3!=128D,\qquad d\le16384D^2.
\]
The two eigenvalues of $A$ on $W$ are
$(2\pm i)^6=-117\pm44i$. For $y_1=y$, $y_2=Ay$, the positive
middle pairing thus has Gram matrix
\[
 d\begin{pmatrix}1&-117\\-117&15625\end{pmatrix},
 \qquad 15625-117^2=1936.
\]
Write $b_j=\langle\delta t,y_j\rangle\in\mathbf Z$. Direct
matrix inversion, cancelling the common scalar before assembling
cycles, gives the integral identity
\[
 1936\delta d\,t=(15625b_1+117b_2)y_1+(117b_1+b_2)y_2.
\]
The coefficient-dual vectors $15625y_1+117y_2$ and
$117y_1+y_2$ have norms $5500\sqrt d$ and $44\sqrt d$.
The signed parts of $y_2$ have degrees at most $125\cdot400D$.
Cauchy--Schwarz therefore bounds the degree of either effective
sign by
\[
 400D\,\delta\sqrt{Td}\,(5500+125\cdot44)
 \le563\,200\,000\delta\sqrt T\,D^2.
\]
The multiplier is at most
$1936\cdot16384\delta D^2=31\,719\,424\delta D^2$.
The stated ceiling produces an integer bound. No coefficient in
the cycle group is divided.
\end{proof}

\begin{theorem}[SF-07: an explicit Eisenstein sixfold division constant]
\label{SF-07}
For a polarized Eisenstein Weil sixfold with integral
$\mathbf Z[\omega]$ action, $\omega^2+\omega+1=0$, a detector of
degree at most $D$ gives
\[
 c_t\le\delta D^2\max\{70\,858\,800,
          270\,050\,760\lceil\sqrt{3T}\rceil\},
 \qquad \delta t\text{ integral},\quad T=-\int t^2>0.
\]
This applies in particular to the $(1^3,3^3)$-polarized cubic source
with $D=D_{E,6}$; it does not relabel that polarization as principal.
\end{theorem}
\begin{proof}
Put $I=[\omega]^*$ and $S=[2+\omega]^*$ on $H^6$. The sum
$1+I+I^2$ kills all character indices except $j=0,3,6$.
The $S$-eigenvalues there are $-27,27,-27$. Hence
\[
 (1+I+I^2)(27-S)=162e_W.
\]
Its positive part on $[Y]$ has degree $3\cdot27D=81D$, and
its negative part has the same bound since $\operatorname{Nm}(2+\omega)=3$.
For $y=162e_W[Y]$ the Theorem~W estimates give
\[
 d=-\int y^2>0,\qquad \sqrt d\le27D,\qquad d\le729D^2.
\]
Take $A=[3+\omega]^*$. Its norm degree factor is $7^3=343$;
its Weil eigenvalues are
\[
 (3+\omega)^6=37+360\omega=-143+180\sqrt{-3}
\]
and its conjugate. The Gram matrix of $y_1=y,y_2=Ay$ is
\[
 d\begin{pmatrix}1&-143\\-143&117649\end{pmatrix},
 \qquad117649-143^2=97200.
\]
For $b_j=\langle\delta t,y_j\rangle$ the exact integral equation is
\[
 97200\delta d\,t=(117649b_1+143b_2)y_1+(143b_1+b_2)y_2.
\]
The coefficient-dual norms are $61740\sqrt{3d}$ and
$180\sqrt{3d}$. Signed representatives of $y_1,y_2$ have each
sign of degree at most $81D$ and $343\cdot81D$, respectively.
Thus either sign in the displayed identity has degree at most
\[
 81D\,\delta\sqrt{3Td}\,(61740+343\cdot180)
 \le270\,050\,760\delta\sqrt{3T}\,D^2.
\]
The multiplier is bounded by
$97200\cdot729\delta D^2=70\,858\,800\delta D^2$.
All coefficients are integral and the claimed ceiling is safe.
\end{proof}
All seven character eigenvalues, both projector identities, both
Gram determinants, the four dual-vector norms and the displayed
integer constants were checked by exact algebra in
\texttt{sixfold\_division\_certificates.py}. This finite computation
does not purport to verify the geometric existence of the detectors
or the supplied Theorem~W. With a smaller integral order one instead
uses HD-WQ with endomorphisms lying in that order and its explicitly
charged projector.

\subsection{Interpretation and the remaining scope}
The new detector bounds give the uniform integral form of condition A,
cote~16 scan~050, with the correction SIGN\_NEXT SN-C03/C04. Their
positivity/division step is condition I, cote~16 scan~049 and cote~44
Theorem~4. Positivity is applied only after the effective supplier
exists; it supplies no cycle in an absent component. The integral
projectors implement the Cayley--Hamilton inverse mechanism of
cote~16 scans~048--049.

The endomorphisms and character projectors are algebraic
correspondences; the resulting cycle has nonzero component in the
motive whose Betti realization is $\bigwedge_K^6H^1$. The covers,
polarizations, finite-type cycle spaces and their maps spread over a
finitely generated integral $\mathbf Z$-algebra after inverting the
finite cover, polarization and level denominators. Their degree-six
classes have weight six and Tate twist $(3)$ in etale realization.
This does not replace uniform cycle degree by pointwise algebraicity
over finite fields or by an $L$-value estimate. The coherent
$K_0$-classes of the detecting subvarieties recover the cycle classes
through the topological filtration and Chern character; this
interpretation makes no semiregularity claim.

SF-01 and SF-02 settle the proposed universal cyclic dominance in the
negative. They leave the broader bounded-division question on the
missing families to other constructions, including the independently
proved eightfold contraction where applicable. Merely adding a
compensating surface changes dimension to eight; a sixfold dominance
theorem by itself gives no theorem in that higher dimension. Any
hypothetical failure is a ghost only with the stated 032-dependent
background of GHOSTS B03. No new archive, EGA, SGA or FGA locator is
invented for these interpretations; the paper's separately verified
spread and cycle-space locators apply.

\begin{proposition}[SF-04a: an explicit Gaussian lattice denominator]
\label{SF-04a}
Let $(A,\theta)$ be a Gaussian Weil sixfold with integral maximal-order
$\mathbf Z[i]$ action, split rational Hermitian form, and
\[
 \Delta=\chi(\theta)=\int_A\theta^6/6!.
\]
Then
\[
 \delta_W(A,\theta)\le(2\Delta)^6D_{G,6}.
\]
More generally, if only the order $\mathbf Z+f\mathbf Z[i]$ acts
integrally, the safe explicit bound is
\[
 \delta_W(A,\theta)\le(2f^{12}\Delta)^6D_{G,6}.
\]
The sharper principal bound of SF-05 is retained when its hypotheses
apply.
\end{proposition}
\begin{proof}
Let $\Lambda$ be the maximal-order Hermitian lattice. Choose a
rational maximal isotropic $K$-subspace and intersect it with
$\Lambda$. The intersection is a saturated Gaussian submodule of
rank three. Since $\mathbf Z[i]$ is a principal ideal domain, it is
a direct summand and its basis extends to an integral basis of
$\Lambda$. In this basis the Hermitian matrix is
\[
 H=\begin{pmatrix}0&B\\B^*&C\end{pmatrix},\qquad
 B\in M_3(\mathbf Z[i]),\quad C=C^*.
\]
Here $B$ is invertible over $K$. Put $b=\det B$. The determinant
of the underlying integral alternating polarization is $\Delta^2$,
so the Hermitian determinant is $-\Delta$; hence
$\operatorname{Nm}(b)=\Delta$. Define
\[
 C'=(B^*)^{-1}CB^{-1}
   =\frac{(\operatorname{adj}B)^*C\operatorname{adj}B}{\Delta},
 \qquad
 T=\begin{pmatrix}1&-C'/2\\0&B^{-1}\end{pmatrix}.
\]
Direct multiplication gives
\[
 T^*HT=\begin{pmatrix}0&1\\1&0\end{pmatrix},\qquad
 T^{-1}=\begin{pmatrix}1&(B^*)^{-1}C/2\\0&B\end{pmatrix}.
\]
Since $b\bar b=\Delta$, all entries of both $T$ and $T^{-1}$
belong to $(2\Delta)^{-1}\mathbf Z[i]$. In particular
$(2\Delta)T$ maps the even unimodular Gaussian source lattice
integrally into the target lattice and pulls its polarization back
to $(2\Delta)^2\Xi$. Transport the target complex structure through
$T$ and use the corresponding source detector from SF-03. Its
threefold degree under pushforward is multiplied by at most
$(2\Delta)^6$, proving the first assertion. The size of the basis
coefficients does not affect this denominator estimate.

For the order of conductor $f$, enlarge the lattice to
$\Lambda'=\mathbf Z[i]\Lambda$. Then
$\Lambda\subset\Lambda'\subset f^{-1}\Lambda$ and
$k=[\Lambda':\Lambda]\le f^{12}$. On
$A'=V_{\mathbf R}/\Lambda'$ the form $f^2\theta$ is an integral
ample polarization $\theta'$ compatible with the maximal order,
and
\[
 \Delta'=\chi(\theta')=f^{12}\Delta/k\le f^{12}\Delta.
\]
The map $[f]:A'\to A$ induced by multiplication by $f$ is an
isogeny with $[f]^*\theta=\theta'$. Consequently pushing a
threefold detector forward has no degree cost relative to these
specified polarizations. Apply the first assertion to $\Delta'$
and use its upper bound. Scaling the Hermitian form by $f^2$ does
not change its determinant class, so splitness is preserved.
\end{proof}
